

% 3. Environnement unique et superposition directe
\begin{tikzpicture}
  % On dessine les deux éléments sur le MÊME repère absolu
  %\clip[](-4, -4) rectangle (11 , 4) ;

    \newcommand{\maCourbe}{
            plot[
                smooth,
                tension=0.5
            ] coordinates {
                (-2,1)(2,2)(8,8)(9,10)
            }
    }

    \draw[
        color=colorThree,
        line width=0.8ex
    ] \maCourbe;

    \begin{scope}
        
        \clip[](2, 2) rectangle (8 ,8) ;
            
        \draw[
            shift = {(0, 0)},
            color=colorOne,
            line width=0.8ex
        ] \maCourbe;

    \end{scope}

    \draw
        (-1.3,0)edge[thick,line width=0.9ex,->,color=\colorOne]node[pos=0.9,below,font=\sffamily\bfseries\fontsize{16pt}{18pt}\selectfont,text=\colorOne]{$ $}(10.8,0)

        (0,-1.3)edge[thick,line width=0.9ex,->,color=\colorOne]node[pos=0.85,left,font=\sffamily\bfseries\fontsize{16pt}{18pt}\selectfont,text=\colorOne]{$ $}(0,10.8)

        (-0.5,2)edge[dashed,line width=0.9ex,color=\colorThree]node[pos=0,left,font=\sffamily\bfseries\fontsize{16pt}{18pt}\selectfont,text=\colorThree]{$ $}(10.8,2)
        (-0.5,2)edge[line width=0.9ex,color=\colorOne]node[pos=0,left,font=\sffamily\bfseries\fontsize{16pt}{18pt}\selectfont,text=\colorOne]{$\bm{y(x,t)}$}(0.5,2)

        (-0.5,8)edge[dashed,line width=0.9ex,color=\colorThree]node[pos=0,left,font=\sffamily\bfseries\fontsize{16pt}{18pt}\selectfont,text=\colorThree]{$\bm{y(x+\mathrm{d}x,t)}$}(10.8,8)
        (-0.5,8)edge[line width=0.9ex,color=\colorOne]node[pos=0,left,font=\sffamily\bfseries\fontsize{16pt}{18pt}\selectfont,text=\colorOne]{$\bm{y(x+\mathrm{d}x,t)}$}(0.5,8)

        (2,-0.5)edge[dashed,line width=0.9ex,color=\colorThree]node[pos=0,below,font=\sffamily\bfseries\fontsize{16pt}{18pt}\selectfont,text=\colorThree]{$ $}(2,10.8)
        (2,-0.5)edge[line width=0.9ex,color=\colorOne]node[pos=0,below,font=\sffamily\bfseries\fontsize{16pt}{18pt}\selectfont,text=\colorOne]{$\bm{x}$}(2,0.5)

        (8,-0.5)edge[dashed,line width=0.9ex,color=\colorThree]node[pos=0,below,font=\sffamily\bfseries\fontsize{16pt}{18pt}\selectfont,text=\colorThree]{$ $}(8,10.8)
        (8,-0.5)edge[line width=0.9ex,color=\colorOne]node[shift = {(0,0.1)}, pos=0,below,font=\sffamily\bfseries\fontsize{16pt}{18pt}\selectfont,text=\colorOne]{$\bm{x + \mathrm{d} x }$}(8,0.5)

        (2,-1.5)edge[thick,line width=0.9ex,<-,color=\colorOne](4.5,-1.5)
        (5.5,-1.5)edge[thick,line width=0.9ex,->,color=\colorOne](8,-1.5)

        (10.8,2)edge[thick,line width=0.9ex,<-,color=\colorOne](10.8,2.3)
        (10.8,7.7)edge[thick,line width=0.9ex,->,color=\colorOne](10.8,8)

        (2,2)edge[thick,line width=0.9ex,->,color=\colorSix]node[shift = {(0,0)}, pos=0.6,above,font=\sffamily\bfseries\fontsize{16pt}{18pt}\selectfont,text=\colorSix ,rotate = 45]{$\bm{\vec{T}(x,t)}$}(-1,-1)

        (8,8)edge[thick,line width=0.9ex,->,color=\colorSix]node[shift = {(0,0)}, pos=0.5,above,font=\sffamily\bfseries\fontsize{16pt}{18pt}\selectfont,text=\colorSix ,rotate = 55]{$\bm{\vec{T}(x + \mathrm{d} x ,t)}$}(11,12)

        (5.5,5)edge[thick,line width=0.9ex,->,color=\colorSix]node[shift = {(0,0)}, pos=0.5,below,font=\sffamily\bfseries\fontsize{13pt}{18pt}\selectfont,text=\colorSix ,rotate = 90]{$\bm{d\vec{P} = - g \, \mathrm{d} m \, \vec{e}_y }$}(5.5,2.2)


    ;

    \node[font=\sffamily\bfseries\fontsize{16pt}{18pt}\selectfont,text=\colorOne] at (5,-1.5) {$\bm{\mathrm{d} x}$};

    \node[font=\sffamily\bfseries\fontsize{16pt}{18pt}\selectfont,text=\colorOne , rotate = 90 ] at (10.8,5) {$\bm{\mathrm{d}y(x,t) + o\left ( \left \vert \frac{\mathrm{d} x}{x} \right \vert ^2 \right )}$};


    \draw[color=colorFour, line width=0.9ex , shift = {(8,8)}] 
        (0,0) -- 
        (1,0) arc (0:53:1cm) 
        node[midway, right=2pt , font=\sffamily\bfseries\fontsize{16pt}{18pt}\selectfont,text=\colorFour] {$\bm{\alpha( x + \mathrm{d} x , t )}$}
        -- (0,0)
    ;

    \draw[color=colorFour, line width=0.9ex , shift = {(2,2)}] 
        (0,0) -- 
        (1,0) arc (0:45:1cm) 
        node[midway, right=2pt , font=\sffamily\bfseries\fontsize{16pt}{18pt}\selectfont,text=\colorFour] {$\bm{\alpha( x , t )}$}
        -- (0,0)
    ;



%   \node[transform canvas={scale=0.5}] at (15,0) {
%     \begin{tikzpicture}

%         \draw[
%             opacity = 0.5, 
%             fill=colorOne!10!colorTwo,
%             draw=colorOne,
%             line width=2.2ex
%             ] (0,0) ellipse [x radius=5cm, y radius=2cm];

%         \draw[
%             fill=none,
%             draw=colorOne,
%             line width=2.2ex
%             ] (0,1) ellipse [x radius=5cm, y radius=2cm];

%         \draw[
%             color=colorOne,
%             line width=2.2ex,
%             pattern={Lines[
%                         angle=90,
%                         distance=3ex,
%                         line width=1.8ex
%                     ]},
%             pattern color=colorOne!50!colorTwo ,
%             ] (5,0) arc[
%             start angle=0,
%             end angle=-180,
%             x radius=5cm,
%             y radius=2cm
%             ]--(-5,1)
%             (-5,1) arc[
%             start angle=180,
%             end angle=360,
%             x radius=5cm,
%             y radius=2cm
%             ]--(5,0)
%             (5,0) arc[
%             start angle=0,
%             end angle=180,
%             x radius=5cm,
%             y radius=2cm
%             ]--(-5,1)
%             (-5,1) arc[
%             start angle=180,
%             end angle=0,
%             x radius=5cm,
%             y radius=2cm
%             ]--(5,0)
%         ;

%         \draw[
%             color=colorFour,
%             line width=1.0ex,
%             ->
%             ] (0,0) -- ({5*cos(45)},{2*sin(45)})
%             node[pos = 0.5, above , font=\fontsize{20pt}{1pt}\selectfont\bfseries, text=colorFour  , rotate = 30 ]{$r$}
%         ;

%         \draw[
%             fill=colorFour,
%             color=colorFour,
%             line width=0.5ex,
%             ] (0,0) circle (0.1)
%         ;

%          \draw[
%             color=colorFour,
%             line width=0.5ex,
%             >-
%             ] (-6,0) -- (-6,0.2)
%         ;
%         \draw[
%             color=colorFour,
%             line width=0.5ex,
%             -<
%             ] (-6,0.8) -- (-6,1)
%         ;
%         \node[font=\fontsize{10pt}{1pt}\selectfont\bfseries, text=colorFour  , rotate = 90 ] at (-6, 0.5) {$\mathrm{d} r$}; 

%         %\node[font=\fontsize{40pt}{1pt}\selectfont\bfseries, text=colorFour  , rotate = 0 ] at (0, -3) {$\mathrm{d} r \not \ll r $};

%         \begin{scope}[shift={(-4,4)}]
%                 \draw[,
%                     line width=1.5ex, color=colorOne, 
%                 fill=colorSix!20,
%                     pattern={Lines[
%                         angle=90,
%                         distance=3ex,
%                         line width=1.8ex
%                     ]},
%                     pattern color=colorOne!50!colorTwo ]  
%                 (0,0) rectangle (3,2);
%                 \node[font=\fontsize{40pt}{1pt}\selectfont\bfseries, text=colorOne , right ] 
%                     at (3.5,1) {$=$};

%                 \node[font=\fontsize{40pt}{1pt}\selectfont\bfseries, text=colorOne , right ] 
%                     at (5.5,1) {$\mathrm{d}A(r)$};


%         \end{scope}

%     \end{tikzpicture}   
%   };

%   \node[transform canvas={scale=0.5}] at (0,0) {
%     \begin{tikzpicture}

%         \begin{scope}

%             \draw[
%             color=colorSix,
%             line width=2.2ex,
%             draw = none, 
%             pattern={Lines[
%                         angle=45,
%                         distance=3ex,
%                         line width=1.8ex
%                     ]},
%             pattern color=colorSix ,
%             ] (5,0) -- (6,0) 
%             (6,0) arc[
%             start angle=0,
%             end angle=360,
%             x radius=6cm,
%             y radius=6cm
%             ]--(5,0)
%             (5,0) arc[
%             start angle=0,
%             end angle=-360,
%             x radius=5cm,
%             y radius=5cm
%             ]
%         ;

%         \end{scope}

%         \draw[
%             fill=none,
%             draw=colorOne,
%             line width=2.2ex
%             ] (0,0) circle (5);

%         \draw[
%             fill=none,
%             draw=colorOne,
%             line width=2.2ex
%             ] (0,0) circle (6);
        
       

%         \draw[
%             color=colorFour,
%             line width=1.0ex,
%             ->
%             ] (0,0) -- ({5*cos(45)},{5*sin(45)})
%             node[pos = 0.5, above , font=\fontsize{20pt}{1pt}\selectfont\bfseries, text=colorFour  , rotate = 45 ]{$r$}
%         ;

%         \draw[
%             fill=colorFour,
%             color=colorFour,
%             line width=0.5ex,
%             ] (0,0) circle (0.1)
%         ;

%          \draw[
%             color=colorFour,
%             line width=0.5ex,
%             >-
%             ] (-6,0) -- (-5.8,0.)
%         ;
%         \draw[
%             color=colorFour,
%             line width=0.5ex,
%             -<
%             ] (-5.2,0) -- (-5,0)
%         ;
%         \node[font=\fontsize{10pt}{1pt}\selectfont\bfseries, text=colorFour   , rotate = 0 ] at (-5.5, 0.0) {$\mathrm{d} r$}; 

%         %\node[font=\fontsize{40pt}{1pt}\selectfont\bfseries, text=colorFour  , rotate = 0 ] at (0, -7) {$\mathrm{d} r  \ll r $};

%         \begin{scope}[shift={(-4,6.5)}]
%                 \draw[,
%                     line width=1.5ex, color=colorOne, 
%                 fill=colorSix!20,
%                     pattern={Lines[
%                         angle=45,
%                         distance=3ex,
%                         line width=1.8ex
%                     ]},
%                     pattern color=colorSix ]  
%                 (0,0) rectangle (3,2);
%                 \node[font=\fontsize{40pt}{1pt}\selectfont\bfseries, text=colorOne , right ] 
%                     at (3.5,1) {$=$};

%                 \node[font=\fontsize{40pt}{1pt}\selectfont\bfseries, text=colorOne , right ] 
%                     at (5.5,1) {$\Delta A(r)$};


%         \end{scope}

%     \end{tikzpicture}   
%   };


  

  % Grille de repérage
  %\draw[line width=2.5ex , color = red] (-8, -5.5) rectangle (8 , 4.5) ;
  %\drawgrid{-3}{11}{-4}{4}{1}{1}
\end{tikzpicture}